\documentclass[../../main.tex]{subfiles}
\begin{document}
% SOURCE p=235,236; printed=220-221; object=ch07-ex15; number=15; region=bottom-upper
\begin{exercise}\label{exercise:ch07-ordinal-measure}
令 $\Omega$ 为所有可数序数的集合，$\omega_1$ 为第一个不可数序数，$\Omega^*=\Omega\cup\{\omega_1\}$，赋予 $\Omega^*$ 顺序拓扑（见\cref{exercise:ch04-order-topology}）。
\begin{enumerate}[label=\alph*.,ref=(\alph*)]
\item\label{item:ch07-ordinal-compact} $\Omega^*$ 是紧 Hausdorff 空间。（提示：$\Omega^*$ 中不存在无限严格递减序列。）
\item\label{item:ch07-ordinal-not-sigma-compact} $\Omega$ 是 $\Omega^*$ 中的开集，但不 $\sigma$ 紧。
\item\label{item:ch07-ordinal-uncountable-cofinal} 子集 $E\subseteq\Omega$ 不可数，当且仅当对每个 $x\in\Omega$，都存在 $y\in E$ 使 $x<y$。
\item\label{item:ch07-ordinal-countable-intersection} 若 $(E_n)$ 是 $\Omega^*$ 中不可数闭集的序列，则 $\bigcap_{n=1}^\infty E_n$ 不可数。（若 $(x_j)$ 是 $\Omega$ 中的递增序列，且每个 $E_n$ 含无穷多个 $x_j$，则 $\lim_jx_j$ 存在并属于 $\bigcap_nE_n$。）
\item\label{item:ch07-ordinal-borel-dichotomy} 若 $E\in\mathscr B_{\Omega^*}$，则 $E\cup\{\omega_1\}$ 或 $E^c\cup\{\omega_1\}$ 含有不可数闭集。（提示：满足后一条件的所有 $E$ 构成一个 $\sigma$ 代数。）
\item\label{item:ch07-ordinal-nonregular-measure} 在 $\mathscr B_{\Omega^*}$ 上定义 $\mu(E)$：若 $E\cup\{\omega_1\}$ 含不可数闭集，取值一；否则取值零。则 $\mu$ 为测度，$\mu(\{\omega_1\})=0$，但每个包含 $\omega_1$ 的开集 $U$ 均有 $\mu(U)=1$。
\item\label{item:ch07-ordinal-eventual-constant} 若 $f\in C(\Omega^*)$，则存在 $x\in\Omega$，使 $f(y)=f(\omega_1)$ 对所有 $y\geqslant x$ 成立。（原书提示：若 $E_n=\{x:|f(x)-f(\omega_1)|<n^{-1}\}$，则 $E_n$ 可数。）
\item\label{item:ch07-ordinal-associated-point-mass} 对 (f) 中的 $\mu$，泛函 $f\mapsto\int f\,\mathrm d\mu$ 在 $\Omega^*$ 上对应的 Radon 测度是 $\omega_1$ 处的点质量。
\end{enumerate}
\end{exercise}
% SUPPLEMENT: id=v14-ch07-ex15-solution; parent=ch07-ex15; kind=solution+correction+motivation+hypotheses; origin=independent-derivation
\begin{solution}
\textbf{序数空间上的非正则概率测度（重构新增）。}先回顾两个集合论事实：每个 $\gamma<\omega_1$ 的初段 $[0,\gamma]$ 可数；可数个可数序数的上确界仍是可数序数，因为其并是可数个可数集的并。

(a) 若 $\alpha<\beta$ 且两者之间有点 $\gamma$，则开射线 $[0,\gamma)$ 与 $(\gamma,\omega_1]$ 分别包含 $\alpha,\beta$ 且不交；若其间无点，则 $[0,\beta)$ 与 $(\alpha,\omega_1]$ 不交并分别包含它们。因此空间 Hausdorff。

给定任意开覆盖，从 $\beta_0=\omega_1$ 开始。若当前剩余初段的端点 $\beta_j>0$，取覆盖中含 $\beta_j$ 的开集 $U_j$；顺序拓扑保证存在 $\beta_{j+1}<\beta_j$，使 $(\beta_{j+1},\beta_j]\subseteq U_j$。删除已覆盖的这段后，剩余至多为 $[0,\beta_{j+1}]$。这个过程不能无穷进行：若有无限严格递减序列，其值集合有最小元，而下一项更小，矛盾。因此经过有限步到达端点零；再选一个覆盖零的开集，便得到有限子覆盖。空间紧。

(b) $\Omega=[0,\omega_1)$ 是开射线，所以开。若 $K\subseteq\Omega$ 紧，用开初段 $[0,\gamma)$（$\gamma<\omega_1$）覆盖 $K$，可取有限子覆盖，因此 $K$ 有某个可数序数上界。若 $\Omega$ 是可数个这样的紧集的并，把这些上界取上确界仍小于 $\omega_1$，会使整个 $\Omega$ 有可数序数上界，矛盾。故它不 $\sigma$ 紧。

(c) 若 $E$ 可数，其上确界为某个 $\gamma<\omega_1$，取 $x=\gamma$ 即没有 $y\in E$ 满足 $x<y$。反过来，若该共尾条件失败，则存在 $x<\omega_1$ 使 $E\subseteq[0,x]$，后者可数，所以 $E$ 可数。这证明等价性。

(d) 每个 $E_n\cap\Omega$ 仍不可数，由 (c) 在 $\Omega$ 中无界。任给 $\gamma\in\Omega$，按指标次序
\begin{align*}
1;\quad1,2;\quad1,2,3;\quad1,2,3,4;\quad\ldots
\end{align*}
递归选择严格递增的 $x_j>\gamma$，使第 $j$ 项位于相应的 $E_n\cap\Omega$。无界性使每次选择可行。令 $\delta=\sup_jx_j<\omega_1$。这个严格递增列在顺序拓扑中趋于 $\delta$：任意 $\eta<\delta$ 终被某项超过，此后所有项位于 $(\eta,\delta]$。每个 $n$ 在指标表中出现无穷次，其对应子列仍趋于 $\delta$；由 $E_n$ 闭，$\delta\in E_n$。所以交集在 $\Omega$ 内含一个大于任意给定 $\gamma$ 的点，按 (c) 不可数。

(e) 记 $P(E)$ 表示 $E\cup\{\omega_1\}$ 含不可数闭集。$P(E)$ 与 $P(E^c)$ 不能同时成立：否则两个相应不可数闭集按 (d) 的有限情形有不可数交集，但交集包含于 $\{\omega_1\}$，矛盾。令 $\mathscr A$ 为满足 $P(E)$ 或 $P(E^c)$ 的所有子集 $E$。它含空集并对取补封闭。若 $E_n\in\mathscr A$，其中某个 $P(E_n)$ 成立，则 $P(\bigcup_nE_n)$ 成立；若无一成立，则每个 $E_n^c\cup\{\omega_1\}$ 含不可数闭集 $F_n$。由 (d)，$\bigcap_nF_n$ 是不可数闭集，包含于 $(\bigcup_nE_n)^c\cup\{\omega_1\}$，所以并集仍属于 $\mathscr A$。因此 $\mathscr A$ 是 $\sigma$ 代数。

它包含所有闭集 $F$：若 $F$ 不可数，$P(F)$ 直接成立；若 $F$ 可数，$F\cap\Omega$ 有可数序数上界，选更大的 $\gamma<\omega_1$，则不可数闭尾段 $[\gamma,\omega_1]\subseteq F^c\cup\{\omega_1\}$，故 $P(F^c)$ 成立。因此 $\mathscr B_{\Omega^*}\subseteq\mathscr A$，证明所求二分性；对 Borel 集恰有一个分支成立。

(f) 显然 $\mu(\varnothing)=0$。设 Borel 集 $E_n$ 两两不交。至多有一个 $\mu(E_n)=1$：否则两个相应不可数闭集的交集只能在 $\{\omega_1\}$ 中，违反 (d)。若存在一个值一的项，则其闭集见证也见证 $\mu(\bigcup_nE_n)=1$，故可数可加。若全部为零，由 (e) 为每个 $E_n^c\cup\{\omega_1\}$ 选不可数闭集；其交集由 (d) 仍不可数，见证 $P((\bigcup_nE_n)^c)$，从而并集测度零。可数可加性在两种情形都成立。单点 $\{\omega_1\}$ 不含不可数闭集，故其测度零。任何含 $\omega_1$ 的开集 $U$ 都包含一个尾段 $(\gamma,\omega_1]$，因而包含不可数闭集 $[\gamma+1,\omega_1]$，故 $\mu(U)=1$。这证明外正则性在单点 $\{\omega_1\}$ 上失败，$\mu$ 不是 Radon 测度，尽管总质量为一。

(g) \textbf{原提示勘误。}原提示中的严格不等号 $<$ 方向错误：由 $f$ 在 $\omega_1$ 处连续，每个原 $E_n$ 都包含某个尾段，因而不可数，而不是可数。应改为
\begin{align*}
F_n=\{x\in\Omega^*:|f(x)-f(\omega_1)|\geqslant1/n\}.
\end{align*}
$F_n$ 闭且不含 $\omega_1$，由 (a) 紧，由 (b) 的紧子集有界性，$F_n\subseteq\Omega$ 有可数序数上界，因而可数。其可数并也有可数序数上界 $\gamma$。取 $x=\gamma+1<\omega_1$；若 $y\geqslant x$，则 $y$ 不属于任何 $F_n$，所以 $|f(y)-f(\omega_1)|<1/n$ 对所有 $n$ 成立，故 $f(y)=f(\omega_1)$。若所有 $F_n$ 均空，可直接取 $x=0$。论证对实值或复值连续函数同样成立。

(h) 由 (g)，$f-f(\omega_1)$ 只在某个可数初段上可能非零。任何可数集 $H\subseteq\Omega^*$ 都有 $\mu(H)=0$，因为 $H\cup\{\omega_1\}$ 仍可数。又 $f$ 在紧空间上有界，所以
\begin{align*}
\int f\,\mathrm d\mu=f(\omega_1)\mu(\Omega^*)=f(\omega_1).
\end{align*}
点质量 $\delta_{\omega_1}$ 是 Radon 测度：紧集上至多为一；若 Borel 集不含 $\omega_1$，可用开集 $\Omega^*\setminus\{\omega_1\}$ 零质量外罩，若含该点，任何外罩均质量一；开集若含该点，以紧单点内逼近即可，若不含则质量零。其积分正是 $f(\omega_1)$，故由 Riesz 唯一性，泛函对应的 Radon 测度为 $\delta_{\omega_1}$。
\end{solution}
\begin{remark}
\textbf{动机与逐条件说明（重构新增）。}不可数闭集在这里能检测“尾部”的质量，可数个这样的闭集仍有不可数交集，这个比一般紧性的有限交性质更强的事实保证了零一取值函数的\emph{可数}可加性。可数序数的上确界仍可数是 (b)—(g) 的共同基础；不能用任意不可数族替代其中的可数族。(a) 的全空间紧性并不使每个开子集都 $\sigma$ 紧，(b) 正是\cref{theorem:ch07-automatic-regularity}所需额外条件失败的地方。

(g) 的结论比趋于端点值更强：连续函数在整个某个尾段上\emph{完全恒定}，原因是可数个误差阈值的起始位置仍有一个可数的共同上界。因此 $\mu$ 与 $\delta_{\omega_1}$ 对所有连续函数给出相同积分，却在单点上分别取零与一；这说明 Riesz 定理的 Radon 唯一性条件不可忽略。原错误提示已经原样保留，并以上述反向不等号订正，不能把一个实际不可数的集合当作可数来完成证明。
\end{remark}
% END SUPPLEMENT: id=v14-ch07-ex15-solution
\end{document}
